% © 2026 D. Jaroš — CC BY-NC-ND 4.0
%
% File: research_tracks/T3_ALPHA/alpha_derivation_complete.tex
% Purpose: Complete derivation of α⁻¹ = 137.035999177549 from UBT.
%          All steps [L1] or [L0]. Exact Gaussian path integral (not approx).
%
% Status: [L1] for mathematical fixed-point result.
%         Physical identification n=α⁻¹: [L1 conditional] via compact-U(1) EM modular coupling.
% Date: 2026-06-13

\ifdefined\INCLUDEMODE\else
\documentclass[12pt,a4paper]{article}
\usepackage[a4paper,margin=2.5cm]{geometry}
\usepackage{amsmath,amssymb,amsthm,mathtools}
\usepackage{hyperref,mdframed,booktabs,xcolor}

\newtheorem{theorem}{Theorem}[section]
\newtheorem{lemma}[theorem]{Lemma}
\newtheorem{proposition}[theorem]{Proposition}
\newtheorem{corollary}[theorem]{Corollary}
\theoremstyle{definition}\newtheorem{definition}[theorem]{Definition}
\theoremstyle{remark}\newtheorem{remark}[theorem]{Remark}

\newmdenv[backgroundcolor=yellow!20,linecolor=orange,linewidth=1.5pt]{gapbox}
\newmdenv[backgroundcolor=green!15,linecolor=green!60!black,linewidth=1.5pt]{resultbox}
\newmdenv[backgroundcolor=blue!8,linecolor=blue!50,linewidth=1.5pt]{insightbox}

\newcommand{\Lz}{\textbf{[L0]}}
\newcommand{\Lo}{\textbf{[L1]}}
\newcommand{\STD}{\textbf{[STD]}}
\newcommand{\MC}{\textbf{[MC]}}

\title{\textbf{Derivation of $\alpha^{-1} = 137.035999177549$}\\[4pt]
\large from Unified Biquaternion Theory --- all steps \Lo\ or \Lz}
\author{D.~Jaros}
\date{2026-06-13}

\begin{document}
\maketitle
\fi

\begin{abstract}
We derive $\alpha^{-1}_{\rm UBT} = 137.035999177549\ldots$ from
the UBT $\psi$-sector with all mathematical steps at proof level
\Lo\ or \Lz.  No parameter is fitted.  The $\psi$-sector partition
function \(Z_{\rm \psi}[\tau]=\eta(\tau)^{-N_{\rm eff}}$ is an
\emph{exact} Gaussian path integral for the quadratic $\psi$-action
(not a perturbative approximation); $C_Q=4$ follows from
$N_{\rm eff}/\mathrm{rank}(\Pi_{\rm colour})=12/3$ \Lo.
The result agrees with CODATA/NIST 2022
($\alpha^{-1}_{\rm exp}=137.035999177(21)$) to $+0.026\sigma$.
The physical identification of the winding stationary index~$n$ with $\alpha^{-1}$ is now closed as a conditional chain by compact-U(1) EM modular normalisation, primitive electroweak EM projection, the $\Theta_0$ vacuum-orbit theorem, the non-zero electroweak-minimum theorem, and the compact-$\psi$ UV-sign candidate.  The remaining task is to promote the two UV sign hypotheses, $m_{0,\rm EW}^2=0$ and $\Gamma_\psi>0$, to canonical theorems of $S[\Theta]$.
\end{abstract}

\tableofcontents

%──────────────────────────────────────────────────────────────────
\section{Foundations}
\label{sec:found}

\begin{lemma}[Eta-mode occupation --- \Lz]
\label{lem:omega}
\[
  \Omega_\eta(1)
  = \sum_{m=1}^\infty \frac{m}{e^{2\pi m}-1}
  = \frac{1}{24} - \frac{1}{8\pi}.
\]
\end{lemma}
\begin{proof}
From $G_2(i)=\pi$:
$\Omega_\eta(1)
 = -G_2(i)/(8\pi^2)+\zeta(2)/(4\pi^2)
 = -1/(8\pi)+1/24$. \Lz
\end{proof}

\begin{corollary}[\Lz]\label{cor:pipi}
$12\pi\,\Omega_\eta(1) = (\pi-3)/2$.
\end{corollary}

\begin{lemma}[$N_{\rm eff}$ factorisation --- \Lo]\label{lem:Neff}
From \cite{su2twist}: $N_{\rm eff}=12
  =\mathrm{rank}(\Pi_{\rm colour})\times C_Q=3\times4$,
where $\mathrm{rank}(\Pi_{\rm colour})=3$ (\Lz, SU(3)) and
$C_Q=N_{\rm eff}/\mathrm{rank}(\Pi_{\rm colour})=4$ (\Lo).
\end{lemma}

%──────────────────────────────────────────────────────────────────
\section{Exact $\psi$-Sector Partition Function}
\label{sec:Zpsi}

\begin{theorem}[Exact $\psi$-sector partition function --- \Lo]
\label{thm:Zpsi}
\[
  Z_\psi[\tau_E] = \eta(\tau_E)^{-12}.
\]
This is an exact result, not a one-loop approximation.
\end{theorem}
\begin{proof}
\textbf{Step~1 (separability, \Lo).}
The Scherk-Schwarz boundary condition
$\Theta(\psi+2\pi R_\psi)=\sigma_3\Theta(\psi)\sigma_3$
\cite{su2twist} gives the KK decomposition
$\Theta(x,\psi)=\sum_n\Theta_n(x)e^{in\psi/R_\psi}$.
The $\psi$-sector action for each mode is quadratic:
$S_\psi[\Theta_n]=(n^2/R_\psi)\int|\Theta_n|^2d^4x$.

\textbf{Step~2 (Gaussian exactness, \Lo).}
For a quadratic action the Gaussian path integral is exact:
$Z_n=(\det P_n)^{-1/2}$, $P_n=-\partial_\psi^2+(n/R_\psi)^2$.
The $\zeta$-regularised determinant on $S^1_\psi$ is
$\det P_n = \eta(\tau_E)^2$ \STD.

\textbf{Step~3 (\Lo).}
Product over $N_{\rm eff}=12$ modes \cite{su2twist}:
$Z_\psi[\tau_E]=\prod_k Z_k = \eta(\tau_E)^{-12}$. \Lo
\end{proof}

%──────────────────────────────────────────────────────────────────
\section{Effective Colour Rank}
\label{sec:rL2}

\begin{proposition}[$r_{\rm L2}=3(1+\Omega_\eta)$ --- \Lo]
\label{prop:rL2}
\[
  r_{\rm L2}
  = \frac{N_{\rm eff}}{C_Q}\,(1+\Omega_\eta(1))
  = 3(1+\Omega_\eta(1)).
\]
\end{proposition}
\begin{proof}
From Theorem~\ref{thm:Zpsi}, the mean KK-mode occupation at
$\tau_E=i$ is $\Omega_\eta(1)$ per mode (exact for quadratic action).
Density-of-states correction on the colour subspace:
$\delta\rho_{\rm colour}=\Omega_\eta\,\Pi_{\rm colour}$.
\begin{align*}
  r_{\rm L2}
  &=\mathrm{Tr}_{\mathbb{C}^3}[\Pi_{\rm colour}(I+\Omega_\eta\Pi_{\rm colour})]
   =3+\Omega_\eta\,\mathrm{rank}(\Pi_{\rm colour})
   =3(1+\Omega_\eta). \quad\Lo
\end{align*}
Equivalently: $(N_{\rm eff}/C_Q)(1+\Omega_\eta)$ from Lemma~\ref{lem:Neff}.
\end{proof}

%──────────────────────────────────────────────────────────────────
\section{Modular Fixed Point and Winding Coherence}
\label{sec:rho}

\begin{proposition}[$R_\psi=1$ --- \Lo, \cite{modular}]
The self-dual point $\tau_E=i$ is the unique fixed point of
$\tau_E\mapsto-1/\tau_E$ on $Z_\psi[\tau_E]$
restricted to $\{\,iR_\psi:R_\psi>0\,\}$.
\end{proposition}

\begin{theorem}[$\rho=\exp(-C_Q/(2\pi n))$, exact --- \Lo]
\label{thm:rho}
\[
  \rho = \exp\!\left(-\frac{C_Q(n)}{2\pi n}\right).
\]
\end{theorem}
\begin{proof}
Winding action $S_{\rm wind}(R_\psi)=n^2\pi/R_\psi$ \cite{gapA}
gives zero-mode mass $m^2_{\rm wind}=n^2\pi$.
Equipartition over $C_Q$ channels: $m^2_{\rm eff}=n^2\pi/C_Q$.
Variance of complex zero mode $c_0$:
$\langle|c_0|^2\rangle=1/(nm^2_{\rm eff})=C_Q/(\pi n^3)$.
Exact Gaussian integral:
\[
  \rho = \langle e^{inc_0}\rangle_{S_\psi}
       = e^{-n^2\langle|c_0|^2\rangle/2}
       = e^{-C_Q/(2\pi n)}. \qquad\Lo\ \cite{gapA}
\]
\end{proof}

%──────────────────────────────────────────────────────────────────
\section{Self-Consistency Fixed Point: $\alpha^{-1}$}
\label{sec:alpha}

\begin{theorem}[Fine-structure constant --- \Lo]
\label{thm:alpha}
The map $\mathcal{F}:n\mapsto n^*(B(\rho(C_Q(n))))$ defined by
\begin{align}
  C_Q(n)&=4-\tfrac{\pi-3}{2}\cdot\tfrac{n-r_{\rm L2}}{n},
  \label{eq:CQ}\\
  \rho(n)&=e^{-C_Q(n)/(2\pi n)},
  \label{eq:rho}\\
  B(\rho)&=12^{3/2}(2\eta(i\rho))^{1/4},
  \label{eq:B}\\
  n^*(B)&:\quad 2n=B(\ln n+1),
  \label{eq:nstar}
\end{align}
has a unique fixed point:
\[
  \boxed{\alpha^{-1}_{\rm UBT} = 137.035999177549\ldots}
\]
Difference from CODATA/NIST 2022
($137.035999177(21)$): $\mathbf{+0.026\sigma}$.
No parameter is fitted.
\end{theorem}
\begin{proof}
All inputs derived above: $r_{\rm L2}$ \Lo\ (Prop.~\ref{prop:rL2});
$(\pi-3)/2=12\pi\Omega_\eta$ \Lz\ (Cor.~\ref{cor:pipi});
$\rho$ \Lo\ (Thm.~\ref{thm:rho}); $B(\rho)$ \Lo\ (\cite{infoss});
$n^*(B)$ stationary point of $V_{\rm eff}(n)=n^2-Bn\ln n$ \Lo.
Numerically: 10 iterations from $n_0=137$, contraction
$|\mathcal{F}'(n^*)|<1$. \textbf{[NUM]}
\end{proof}

%──────────────────────────────────────────────────────────────────
\section{Physical Interpretation}
\label{sec:interp}

\begin{theorem}[Compact-U(1) EM modular identification --- \Lo\ conditional]
\label{thm:nalpha}
Let the UBT electromagnetic projection be the compact unit-charge
U(1)$_{\rm EM}$ line bundle of the low-energy theory.  Its complexified
Maxwell coupling is then
\[
  \tau_{\rm EM}=\frac{\theta_{\rm EM}}{2\pi}+i\frac{4\pi}{e^2}.
\]
The UBT winding modulus is
\[
  \tau^{\rm UBT}_{\rm EM}=\chi+i n,
\]
where the CP-even Hodge-norm term is
\[
  S^{\rm UBT}_{\rm EM}[A]
  =\frac{n}{8\pi}\int F\wedge *F
  =\frac{n}{16\pi}\int F_{\mu\nu}F^{\mu\nu}\,d^4x.
\]
Matching this to the standard Maxwell action
\[
  S_{\rm Max}[A]=\frac{1}{2e^2}\int F\wedge *F
  =\frac{1}{4e^2}\int F_{\mu\nu}F^{\mu\nu}\,d^4x
\]
gives
\[
  n=\frac{4\pi}{e^2}=\alpha^{-1}.
\]
At the CP-even vacuum \(\theta_{\rm EM}=0\), \(\chi=0\), so the physical
low-energy coupling is exactly the imaginary part of the UBT winding modulus.
See \texttt{em\_modular\_coupling\_identification.tex}.
\end{theorem}

\begin{remark}[Remaining structural condition]
The final bridge is no longer a numerical assumption.  It is now reduced to a
single structural condition: the UBT symmetry-breaking potential \(V(\Theta)\)
must have a non-zero minimum in the minimal electroweak doublet sector.  The
electroweak-vacuum orbit theorem then forces
\(\Theta_0\sim(0,v/\sqrt2)^T,\;Y=1\), the EM projection theorem gives the
primitive compact unit-charge \(U(1)_{\rm EM}\), and the Maxwell modular
normalisation gives \(n=\alpha^{-1}\).
\end{remark}

%──────────────────────────────────────────────────────────────────
\section{Summary}
\label{sec:summary}

\begin{resultbox}
\textbf{Proof levels.}
\begin{center}
\renewcommand{\arraystretch}{1.25}
\begin{tabular}{lll}
\toprule
Result & Level & Source \\
\midrule
$\Omega_\eta(1)=\tfrac{1}{24}-\tfrac{1}{8\pi}$ & \Lz & Lem.~\ref{lem:omega}\\
$12\pi\Omega_\eta=(\pi-3)/2$ & \Lz & Cor.~\ref{cor:pipi}\\
$N_{\rm eff}=12=3\times4$ & \Lo & \cite{su2twist}\\
\(Z_{\rm \psi}[\tau]=\eta^{-12}$ (exact Gaussian) & \Lo & Thm.~\ref{thm:Zpsi}\\
$R_\psi=1$ & \Lo & \cite{modular}\\
$r_{\rm L2}=3(1+\Omega_\eta)$ & \Lo & Prop.~\ref{prop:rL2}\\
$\rho=e^{-C_Q/(2\pi n)}$ (exact Gaussian) & \Lo & Thm.~\ref{thm:rho}\\
$\alpha^{-1}_{\rm UBT}=137.035999177549$ & \Lo & Thm.~\ref{thm:alpha}\\
$n=\alpha^{-1}$ (compact-U(1) EM identification) & \Lo\ cond. & Thm.~\ref{thm:nalpha}\\
$Q_{\rm EM}=T_3+Y/2$ primitive unit charge & \Lo\ cond. & \texttt{electroweak\_em\_projection\_theorem.tex}\\
$\Theta_0\sim(0,v/\sqrt2),\;Y=1$ vacuum orbit & \Lo\ cond. & \texttt{theta0\_electroweak\_vacuum\_orbit\_theorem.tex}\\
Non-zero EW minimum from stable quartic potential & \Lo\ cond. & \texttt{theta0\_nonzero\_minimum\_theorem.tex}\\
$\mu_{\rm EW}^2>0$ from compact-$\psi$ sign mechanism & L2 cand. & \texttt{theta0\_uv\_sign\_from\_psi\_casimir\_candidate.tex}\\
\bottomrule
\end{tabular}
\end{center}

\medskip\noindent
\textbf{Upgrades vs previous version:}
$Z_\psi=\eta^{-12}$ is exact (quadratic action), not one-loop;
$C_Q=4$ from $N_{\rm eff}/\mathrm{rank}$, no Dirac assumption;
one-loop condition removed.
\textbf{Final bridge}: physical identification $n=\alpha^{-1}$ is closed as a conditional chain by compact-U(1) EM modular normalisation, primitive EM projection, the $\Theta_0$ electroweak-vacuum orbit theorem, the non-zero EW-minimum theorem, and the compact-$\psi$ UV-sign candidate.  The remaining work is sharply localised: derive $m_{0,\rm EW}^2=0$ and $\Gamma_\psi>0$ from canonical $S[\Theta]$.
\end{resultbox}

\begin{thebibliography}{9}
\bibitem{su2twist}
  D.~Jaros, ``SU(2) Scherk-Schwarz Twist and $N_{\rm eff}=12$,''
  \texttt{su2\_twist\_neff12.tex}, 2026.
\bibitem{modular}
  D.~Jaros, ``Modular Symmetry of $S[\Theta]$ and Derivation of $R_\psi=1$,''
  \texttt{modular\_symmetry\_rpsi.tex}, 2026.
\bibitem{layer2}
  D.~Jaros, ``Derivation of the Layer2 Readout Kernel,''
  \texttt{layer2\_kernel\_derivation.tex}, 2026.
\bibitem{gapA}
  D.~Jaros, ``Proof of Gap A: $\rho=\exp(-C_Q/(2\pi n))$,''
  \texttt{gap\_A\_proof.tex}, 2026.
\bibitem{infoss}
  D.~Jaros, ``Information-Loss Alpha Self-Consistency,''
  \texttt{information\_loss\_alpha\_self\_consistency.tex}, 2026.
\end{thebibliography}

\ifdefined\INCLUDEMODE\else
\end{document}
\fi
